\subsection{幂零李代数表示的分解}

设 \(\mathfrak{h}\) 是一个李代数，M 是一个 \(\mathfrak{h}\) -模。对任何映射 \(\lambda\) （从 \(\mathfrak{h}\) 到 \(k\) ），我们在第 1 节中已定义了 M 的向量子空间 \(\mathrm{M}^{\lambda}(\mathfrak{h})\) 与 \(\mathrm{M}_{\lambda}(\mathfrak{h})\) 。特别地，若 \(\mathfrak{g}\) 是一个以 \(\mathfrak{h}\) 为子代数的李代数，且 \(x \in \mathfrak{g}\) ，我们将经常使用记号 \(\mathfrak{g}^{\lambda}(\mathfrak{h})\) 与 \(\mathfrak{g}_{\lambda}(\mathfrak{h})\) ；这时应理解为 \(\mathfrak{h}\) 作用在 \(\mathfrak{g}\) 上，经由伴随表示 ad \(_{\mathfrak{g}}\) 。

命题 8. 设 \(\mathfrak{h}\) 是一个李代数，L、M、N 是 \(\mathfrak{h}\) -模。用 P 表示从 \(\mathfrak{h}\) 到 \(k\) 的映射所成的集合。

\begin{enumerate}
\def\labelenumi{(\roman{enumi})}
\item
  和 \(\sum_{\lambda \in \mathrm{P}}\mathrm{L}^{\lambda}(\mathrm{P})\) 是直和。
\item
  若 \(f: \mathrm{L} \to \mathrm{M}\) 是 \(\mathfrak{h}\) -模的同态，则 \(f(\mathrm{L}^{\lambda}(\mathfrak{h})) \subset \mathrm{M}^{\lambda}(\mathfrak{h})\) 对所有 \(\lambda \in \mathrm{P}\) 成立。
\item
  若 \(f:\mathrm{L}\times \mathrm{M}\to \mathrm{N}\) 是 \(\mathfrak{h}\) -不变的双线性映射，
\end{enumerate}

\[
f (\mathrm{L} ^ {\lambda} (\mathfrak {h}) \times \mathrm{M} ^ {\mu} (\mathfrak {h})) \subset \mathrm{N} ^ {\lambda + \mu} (\mathfrak {h})
\]

对所有 \(\lambda, \mu \in \mathrm{P}\) 成立。

这由命题 2 与命题 3 得出。

命题 9. 设 \(\mathfrak{h}\) 是一个幂零李代数，M 是一个有限维 \(\mathfrak{h}\) -模。用 P 表示从 \(\mathfrak{h}\) 到 \(k\) 的映射所成的集合。

\begin{enumerate}
\def\labelenumi{(\roman{enumi})}
\item
  每个 \(\mathrm{M}^{\lambda}(\mathfrak{h})\) 都是 \(\mathfrak{h}\) -模 \(\mathrm{M}\) 的子模。若 \(x_{M}\) 对所有 \(x\in \mathfrak{h}\) 都可三角化，则 \(\mathrm{M} = \sum_{\lambda \in \mathbb{R}}\mathrm{M}^{\lambda}(\mathfrak{h})\) 。
\item
  若 \(k\) 是无限的，则存在 \(x \in \mathfrak{h}\) 使 \(\mathrm{M}^0(x) = \mathrm{M}^0(\mathfrak{h})\) 。
\item
  若 \(k\) 的特征为 0，且 \(\lambda \in \mathrm{P}\) 使 \(\mathrm{M}^{\lambda}(\mathfrak{h}) \neq 0\) ，则 \(\lambda\) 是 \(\mathfrak{h}\) 上在 \([\mathfrak{h},\mathfrak{h}]\) 上为零的线性形式，且 \(\mathrm{M}_{\lambda}(\mathfrak{h}) \neq 0\) 。
\item
  若 \(f: \mathrm{M} \to \mathrm{N}\) 是有限维 \(\mathfrak{h}\) -模的满射同态，则 \(f(\mathrm{M}^{\lambda}(\mathfrak{h})) = \mathrm{N}^{\lambda}(\mathfrak{h})\) 对所有 \(\lambda \in \mathrm{P}\) 成立。
\item
  若 \(\mathrm{N}\) 是有限维 \(\mathfrak{h}\) -模，\(\mathrm{B}\) 是 \(\mathrm{M} \times \mathrm{N}\) 上在 \(\mathfrak{h}\) 下不变的双线性型，则 \(\mathrm{M}^{\lambda}(\mathfrak{h})\) 与 \(\mathrm{N}^{\mu}(\mathfrak{h})\) 关于
\end{enumerate}

B 正交，如果 \(\lambda + \mu \neq 0\) 。此外，若 B 非退化，则 B 在 \(\mathrm{M}^{\lambda}(\mathfrak{h}) \times \mathrm{N}^{-\lambda}(\mathfrak{h})\) 上的限制对所有 \(\lambda \in \mathrm{P}\) 也非退化。

(i) 由第 1 节的引理 1 与定理 1 得出。(ii) 由第 2 节的命题 7 得出。(iv) 由第 1 节定理 1 的推论 3 得出。我们证明 (iii)。可假定 \(\mathrm{M} = \mathrm{M}^{\lambda}(\mathfrak{h})\) 。于是，对所有 \(x \in \mathfrak{h}\) ，\(\lambda(x) = (\dim \mathrm{M})^{-1}\mathrm{Tr}(x_{\mathrm{M}})\) ；这证明 \(\lambda\) 是线性的（这一点也可由命题 5 得出），且 \(\lambda\) 在 \([\mathfrak{h},\mathfrak{h}]\) 上为零。考虑由下式定义的映射 \(\rho : \mathfrak{h} \to \mathrm{End}_k(\mathrm{M})\) ：

\[
\rho (x) = x _ {\mathrm{M}} - \lambda (x) 1 _ {\mathrm{M}};
\]

由前所述，这是 \(\mathfrak{h}\) 在 M 上的一个表示，且 \(\rho(x)\) 对所有 \(x \in \mathfrak{h}\) 幂零。由 Engel 定理（第一章，§4，第 2 节，定理 1），存在 M 中的 \(m \neq 0\) 使得 \(\rho(x)m = 0\) 对所有 \(x \in \mathfrak{h}\) 成立，故 \(m \in \mathrm{M}_{\lambda}(\mathfrak{h})\) 。

(v) 的第一个断言由第 1 节的命题 2 (ii) 得出。为证第二个断言，鉴于第 1 节的命题 1，可假定 \(k\) 是代数闭的；于是它由第一个断言以及 \(\mathrm{M} = \sum_{\lambda} \mathrm{M}^{\lambda}(\mathfrak{h})\) 、\(\mathrm{N} = \sum_{\lambda} \mathrm{N}^{\lambda}(\mathfrak{h})\) 这一事实得出，参看~(i)。

注. 假定 \(k\) 是特征 2 的完全域。设 \(\mathfrak{h} = \mathfrak{sl}(2, k)\) ，并设 \(M\) 是 \(\mathfrak{h}\) -模 \(k^2\) （对应于 \(\mathfrak{h}\) 的恒等映射）。若 \(x = \begin{pmatrix} a & b \\ c & a \end{pmatrix}\) 是 \(\mathfrak{h}\) 的任意元素，用 \(\lambda(x)\) 表示唯一的 \(\lambda \in k\) ，它使 \(\lambda^2 = a^2 + bc\) 。直接计算表明 \(M = M^\lambda(\mathfrak{h})\) ；另一方面，\(M_\lambda(\mathfrak{h}) = 0\) ，且 \(\lambda\) 既不是线性的也不在 \([\mathfrak{h}, \mathfrak{h}]\) 上为零，尽管 \(\mathfrak{h}\) 是幂零的。

推论. 设 \(\mathfrak{h}\) 是一个幂零李代数，M 是一个有限维 \(\mathfrak{h}\) -模且满足 \(\mathrm{M}^0 (\mathfrak{h}) = 0\) 。设 \(f:\mathfrak{h}\to \mathrm{M}\) 是线性映射，使得

\[
f ([ x, y ]) = x. f (y) - y. f (x) \quad \text { for } x, y \in \mathfrak {h}.
\]

存在 \(a \in \mathbb{M}\) 使得 \(f(x) = x.a\) 对所有 \(x \in \mathfrak{h}\) 成立。

令 \(N = M \times k\) 。让 h 按如下公式作用在 N 上：

\[
x. (m, \lambda) = (x. m - \lambda f (x), 0).
\]

\(f\) 所满足的恒等式表明 \(\mathbf{N}\) 是一个 \(\mathfrak{h}\) -模（第一章，§1，第 8 节，例 2）。映射 \((m, \lambda) \mapsto \lambda\) （从 \(\mathbf{N}\) 到 \(k\) ）是从 \(\mathbf{N}\) 到平凡 \(\mathfrak{h}\) -模 \(k\) 的同态。由命题 9 (iv)，\(\mathrm{N}^0(\mathfrak{h})\) 含有一个形如 \((a, 1)\) 的元素，其中 \(a \in \mathrm{M}\) 。鉴于对 \(\mathbf{M}\) 的假设，

\[
(\mathrm{M} \times 0) \cap \mathrm{N} ^ {0} (\mathfrak {h}) = 0,
\]

故 \(\mathrm{N}^0 (\mathfrak{h})\) 的维数为 1，从而被 \(\mathfrak{h}\) 零化。于是，\(x.a - f(x) = 0\) 对所有 \(x\in \mathfrak{h}\) 成立，这就证明了该推论。

命题 10. 设 \(\mathfrak{g}\) 是一个李代数，\(\mathfrak{h}\) 是 \(\mathfrak{g}\) 的一个幂零子代数。用 P 表示从 \(\mathfrak{h}\) 到 \(k\) 的映射所成的集合。

\begin{enumerate}
\def\labelenumi{(\roman{enumi})}
\item
  对 \(\lambda, \mu \in \mathrm{P}\) ，\([\mathfrak{g}^{\lambda}(\mathfrak{h}), \mathfrak{g}^{\mu}(\mathfrak{h})] \subset \mathfrak{g}^{\lambda + \mu}(\mathfrak{h})\) ；特别地，\(\mathfrak{g}^{0}(\mathfrak{h})\) 是 \(\mathfrak{g}\) 的一个包含 \(\mathfrak{h}\) 的李子代数，且诸 \(\mathfrak{g}^{\lambda}(\mathfrak{h})\) 在 \(\operatorname{ad} \mathfrak{g}^{0}(\mathfrak{h})\) 下稳定。此外，\(\mathfrak{g}^{0}(\mathfrak{h})\) 是它在 \(\mathfrak{g}\) 中的自身的正规化子。
\item
  若 M 是一个 g-模，则 \(\mathfrak{g}^{\lambda}(\mathfrak{h})\mathrm{M}^{\mu}(\mathfrak{h})\subset\mathrm{M}^{\lambda+\mu}(\mathfrak{h})\) 对 \(\lambda,\mu\in P\) 成立；特别地，每个 \(\mathrm{M}^{\lambda}(\mathfrak{h})\) 都是一个 \(\mathfrak{g}^{0}(\mathfrak{h})\) -模。
\item
  若 B 是 g 上在 h 下不变的双线性型，则 \(\mathfrak{g}^{\lambda}(\mathfrak{h})\) 与 \(\mathfrak{g}^{\mu}(\mathfrak{h})\) 关于 B 正交，如果 \(\lambda + \mu \neq 0\) 。假定 B 非退化。那么，对所有 \(\lambda \in P\) ，B 在 \(\mathfrak{g}^{\lambda}(\mathfrak{h}) \times \mathfrak{g}^{-\lambda}(\mathfrak{h})\) 上的限制非退化；特别地，B 在 \(\mathfrak{g}^{0}(\mathfrak{h}) \times \mathfrak{g}^{0}(\mathfrak{h})\) 上的限制非退化。
\item
  假定 \(k\) 的特征为 0。那么，若 \(x \in \mathfrak{g}^{\lambda}(\mathfrak{h})\) 且 \(\lambda \neq 0\) ，则 ad \(x\) 幂零。
\end{enumerate}

由 Jacobi 恒等式，映射 \((x,y)\mapsto [x,y]\) （从 \(\mathfrak{g}\times \mathfrak{g}\) 到 \(\mathfrak{g}\) ）是 \(\mathfrak{g}\) -不变的，从而是 \(\mathfrak{h}\) -不变的。于是 (i) 的第一部分由命题 2 (ii) 得出。(ii) 的证明类似。

若 \(x\) 属于 \(\mathfrak{g}^0(\mathfrak{h})\) 在 \(\mathfrak{g}\) 中的正规化子，则 (ad \(y\) ). \(x = -[x, y] \in \mathfrak{g}^0(\mathfrak{h})\) 对所有 \(y \in \mathfrak{h}\) 成立，故 \((\operatorname{ad} y)^n\) . \(x = 0\) 对充分大的 \(n\) 成立。这证明 \(x \in \mathfrak{g}^0(\mathfrak{h})\) 。断言 (i) 至此完全证得。

断言 (iii) 由命题 9 (v) 得出。

为证 (iv)，可假定 \(k\) 是代数闭的。设 \(x \in \mathfrak{g}^{\lambda}(\mathfrak{h})\) ，其中 \(\lambda \neq 0\) 。对所有 \(\mu \in \mathrm{P}\) 与任何整数 \(n \geq 0\) ，\((\operatorname{ad} x)^n \mathfrak{g}^\mu(\mathfrak{h}) \subset \mathfrak{g}^{\mu + n\lambda}(\mathfrak{h})\) ；设 \(\mathrm{P}_1\) 是 \(\mu \in \mathrm{P}\) 中使 \(\mathfrak{g}^\mu(\mathfrak{h}) \neq 0\) 的元素所成的有限集合；若 \(k\) 的特征为 0 且 \(\lambda \neq 0\) ，则 \((\mathrm{P}_1 + n\lambda) \cap \mathrm{P}_1 = \varnothing\) 对充分大的 \(n\) 成立，故 \((\operatorname{ad} x)^n = 0\) 。

引理 2. 假定 k 的特征为 0。设 g 是 k 上的一个半单李代数，B 是 g 的 Killing 型，m 是 g 的一个子代数。假定下列条件满足：

\begin{enumerate}
\def\labelenumi{\arabic{enumi})}
\item
  B 在 m 上的限制非退化；
\item
  若 \(x \in \mathfrak{m}\) ，则半单分量与幂零分量\footnote{由第一章，§6，第 3 节，定理 3，每个 \(x \in \mathfrak{g}\) 都可唯一地写成一个半单元素 \(s\) 与一个幂零元素 \(n\) 之和，且二者相互交换；元素 \(s\) （相应地，\(n\) ）称为 \(x\) 的半单（相应地，幂零）分量。}（\(x\) 在 \(\mathfrak{g}\) 中的）都属于 \(\mathfrak{m}\) 。
\end{enumerate}

那么 \(\mathfrak{m}\) 在 \(\mathfrak{g}\) 中是约化的（第一章，§ 6，第 6 节）。

由第一章，§6，第 4 节，命题 5 d)，m 是约化的。设 c 是 m 的中心。若 \(x \in c\) 幂零，则 x = 0；事实上，对所有 \(y \in m\) ，ad x 与 ad y 交换，它们的复合 ad \(x \circ ad y\) 幂零，且 \(\mathrm{B}(x, y) = 0\) ，故 x = 0。现在设 x 是 c 的任意元素；设 s 与 n 是它的半单分量与幂零分量。我们有 \(n \in m\) 。由于 ad n 形如 \(\mathrm{P}(\mathrm{ad}x)\) ，其中 P 是无常数项的多项式，故 (ad n).m = 0，从而 \(n \in c\) ，再由前述即得 n = 0。于是 ad x 是半单的。因此，g 的伴随表示在 m 上的限制是半单的（第一章，§6，第 5 节，定理 4 b)）。

命题 11. 假定 k 的特征为 0。设 g 是一个半单李代数，h 是 g 的一个幂零子代数。代数 \(\mathfrak{g}^{0}(\mathfrak{h})\) 满足引理 2 的条件 (1) 与 (2)；它在 g 中是约化的。

设 \(x, x' \in \mathfrak{g}\) ，\(s\) 与 \(s'\) 是它们的半单分量，\(n\) 与 \(n'\) 是它们的幂零分量。我们有

\[
\begin{array}{r l} x ^ {\prime} \in \mathfrak {g} ^ {0} (x) & \Longleftrightarrow (\operatorname{ad} s) (x ^ {\prime}) = 0 \quad \text {(Prop. 4)} \\ & \Longleftrightarrow (\operatorname{ad} x ^ {\prime}) (s) = 0 \\ & \Longrightarrow (\operatorname{ad} s ^ {\prime}) (s) = 0 \\ & \Longleftrightarrow (\operatorname{ad} s) (s ^ {\prime}) = 0 \\ & \Longleftrightarrow s ^ {\prime} \in \mathfrak {g} ^ {0} (x) \quad \text {(Prop. 4)} \end{array}
\]

于是 \(x' \in \mathfrak{g}^{0}(x) \Rightarrow n' \in \mathfrak{g}^{0}(x)\) ，(2) 得证。g 的 Killing 型非退化，故它在 \(\mathfrak{g}^{0}(\mathfrak{h})\) 上的限制非退化（命题 10 (iii)）。于是 \(\mathfrak{g}^{0}(\mathfrak{h})\) 在 g 中约化这一事实由引理 2 得出。

